\let\negmedspace\undefined
\let\negthickspace\undefined
\documentclass[journal]{IEEEtran}
\usepackage[a5paper, margin=10mm, onecolumn]{geometry}
%\usepackage{lmodern} % Ensure lmodern is loaded for pdflatex
\usepackage{tfrupee} % Include tfrupee package

\setlength{\headheight}{1cm} % Set the height of the header box
\setlength{\headsep}{0mm}     % Set the distance between the header box and the top of the text

\usepackage{gvv-book}
\usepackage{gvv}
\usepackage{cite}
\usepackage{amsmath,amssymb,amsfonts,amsthm}
\usepackage{algorithmic}
\usepackage{graphicx}
\usepackage{textcomp}
\usepackage{xcolor}
\usepackage{txfonts}
\usepackage{listings}
\usepackage{enumitem}
\usepackage{mathtools}
\usepackage{gensymb}
\usepackage{comment}
\usepackage[breaklinks=true]{hyperref}
\usepackage{tkz-euclide} 
\usepackage{listings}
% \usepackage{gvv}                                        
\def\inputGnumericTable{}                                 
\usepackage[latin1]{inputenc}                                
\usepackage{color}                                            
\usepackage{array}                                            
\usepackage{longtable}                                       
\usepackage{calc}                                             
\usepackage{multirow}                                         
\usepackage{hhline}                                           
\usepackage{ifthen}
\usepackage{lscape}
\begin{document}

\bibliographystyle{IEEEtran}



\title{5.13.82}
\author{EE25BTECH11057 - Rushil Shanmukha Srinivas
}
% \maketitle
% \newpage
% \bigskip
{\let\newpage\relax\maketitle}

\renewcommand{\thefigure}{\theenumi}
\renewcommand{\thetable}{\theenumi}
\setlength{\intextsep}{10pt} % Space between text and floats

\numberwithin{equation}{enumi}
\numberwithin{figure}{enumi}
\renewcommand{\thetable}{\theenumi}

\textbf{Question :}
 Let $\alpha$ and $\beta$ be the distinct roots of the equation
$
x^2 + x - 1 = 0.
$
Consider the set $T = \{1, \alpha, \beta\}$. For a $3 \times 3$ matrix $\vec{M} = (a_{ij})_{3 \times 3}$, define
$
R_i = a_{i1} + a_{i2} + a_{i3}, 
\newline
\quad C_j = a_{1j} + a_{2j} + a_{3j}, \quad \text{for } i,j = 1,2,3.
$
Match the following:

\begin{enumerate}
\item[(A)] The number of matrices $\vec{M} = (a_{ij})_{3 \times 3}$ with all entries in $T$ such that $R_i = C_j = 0$ for all $i,j$, is


\item[(B)] The number of symmetric matrices $\vec{M} = (a_{ij})_{3 \times 3}$ with all entries in $T$ such that $C_j = 0$ for all $j$, is


\item[(C)] Let $\vec{M} = (a_{ij})_{3 \times 3}$ be a matrix with all entries in $T$ such that $R_i = 0$  for all $i$. Then the absolute value of the determinant of $\vec{M}$ is

\end{enumerate}
\begin{multicols}{3}
\begin{enumerate}
\item 1 \qquad
\item 12
\item infinite
\item 6
\item 0
\end{enumerate}
\end{multicols}
 The correct option is:

\begin{multicols}{2}
\begin{enumerate}
\item $A \to 4, B \to 2, C \to 1$ 
\item $A \to 2, B \to 4, C \to 5$
\item $A \to 2, B \to 4, C \to 3$ 
\item $A \to 1, B \to 5, C \to 4$
\end{enumerate}
\end{multicols}

\textbf{Solution :}

Let $T = \{1,\alpha,\beta\}$ where $\alpha,\beta$ are distinct roots of $x^2+x-1=0$.  
For a $3 \times 3$ matrix $M = (a_{ij})$, define
\begin{align}
R_i = a_{i1} + a_{i2} + a_{i3}, 
\quad 
C_j = a_{1j} + a_{2j} + a_{3j}.
\end{align}
\begin{align}
\alpha+\beta = \frac{-coefficient \ of \ x}{coefficient \ of \ x^2} = -1
\end{align}
\begin{align}
1+\alpha+\beta = 0
\end{align}
\textbf{(A)} We need to find matrices $\vec{M}$ with all entries in $T$ such that $R_i=C_j=0$ for all $i,j$ ,
So
\begin{align}
R_1 = R_2 = R_3 = C_1 = C_2 = C_3 = 0.
\end{align}
Only possible case is some permutation of $\{1,\alpha,\beta\}$ as $1 + \alpha + \beta = 0$ .

As both row total = column total = 0 .

\textbf{Latin Square :}

Latin Square is a Matrix in which one element occurs exactly once in a row and a column.

The Matrix $\vec{M}$ should be Latin Square of order 3 on the symbol set $\{1,\alpha,\beta\}$ 

By permutations , we get
\begin{align}
3\times2\times1\times2 = 12
\end{align}
One of the 12 Latin Square Matrices is given as 
\begin{align}
\vec{M} = \myvec{1 & \alpha & \beta \\
       \alpha & \beta & 1 \\
       \beta & 1 & \alpha}
\end{align}
The number of Matrices $\vec{M}$ = 12.

Thus, A = 2.

\textbf{(B)} Symmetric matrices $\vec{M}$ with entries in $T$ such that $C_j=0$.  
\begin{align}
\vec{M} = \vec{M^\top}
\end{align}
So both $R_i = C_j = 0$ for all i,j. 

\textbf{Latin Square :}

So The Matrix $\vec{M}$ should be Latin Square of order 3 on the symbol set $\{1,\alpha,\beta\}$ and satisfying $\vec{M}=\vec{M}^\top$.

By permutations, we get
\begin{align}
3\times2\times1 = 6
\end{align}
One of the 6 Latin Square Matrices is given as 
\begin{align}
\vec{M} = \myvec{1 & \alpha & \beta \\
       \alpha & \beta & 1 \\
       \beta & 1 & \alpha}
\end{align}
The number of Matrices $\vec{M}$ = 6.

Thus, B = 4.

\textbf{(C)} Matrices $\vec{M}$ with entries in $T$ such that $R_i=0$ for all i. 
So each row is a permutation of $(1,\alpha,\beta)$.  
\begin{align}
1 + \alpha + \beta = 0
\end{align}
If rows are identical
\begin{align}
\mydet{\vec{M}} = 0
\end{align}
If rows are distinct then
For matrix $\vec{M}$ if $C_1 \longrightarrow C_1 + C_2 + C_3$ then every element of $C_1$ is 0.  

So $\mydet{\vec{M}} = 0$.

Thus $C=5$.
\begin{align}
\boxed{A=2, \; B=4, \; C=5}
\end{align}

Correct option is (b).

\end{document}
